\documentclass[11pt,a4paper]{article}

\usepackage[margin=1.8cm]{geometry}
\usepackage{fontspec}
\usepackage{enumitem}
\usepackage{xurl}
\usepackage[hidelinks]{hyperref}
\usepackage{xcolor}
\usepackage{titlesec}

\setmainfont{Times New Roman}
\definecolor{heading}{HTML}{003366}
\definecolor{subtle}{HTML}{555555}

\pagestyle{empty}
\setlength{\parindent}{0pt}
\setlength{\parskip}{5pt}
\sloppy
\setlist[itemize]{leftmargin=1.35em, topsep=2pt, itemsep=2pt}

\titleformat{\section}{\large\bfseries\color{heading}}{}{0pt}{}[\titlerule]
\titleformat{\subsection}{\normalsize\bfseries\color{heading}}{}{0pt}{}

\begin{document}

{\LARGE \textbf{Teaching Statement}}\\[-1pt]
\textbf{Wen Quan}\\
\textcolor{subtle}{Teaching Areas: Interior Design, Decorative Materials, Architectural Design, CAD/BIM, Design Technology, Age-Friendly Design}\\
\href{mailto:wenquan@student.usm.my}{wenquan@student.usm.my} \,|\, 
\href{https://wenquan0816.github.io/schoolar-site/}{\nolinkurl{wenquan0816.github.io/schoolar-site}}

\section*{Teaching Philosophy}
My teaching philosophy is grounded in three principles: design must be evidence-based, technical skills must support human needs, and students learn best when they connect studio thinking with real social problems. In interior design and built environment education, I want students to move beyond attractive visual outcomes and learn how to ask: Who is the space for? What problem does the design solve? How can the design be evaluated?

My teaching experience includes undergraduate courses with typical class sizes of approximately 40 students. I have taught and supported courses related to interior design, decorative materials, architectural design, engineering drawing, architectural surveying, building structures, CAD, BIM software, design software, renovation design, rendering, architectural space design, sketching, architectural investigation, and architectural history. These areas allow me to contribute to both studio-based and technology-based teaching in Malaysian university programs.

\section*{Approach to Studio and Design Teaching}
In studio teaching, I emphasize a structured design process: problem definition, user analysis, site/context investigation, precedent study, concept development, technical resolution, representation, critique, and reflection. I encourage students to treat drawings, models, renderings, and diagrams as tools for reasoning, not only presentation.

For interior design and architectural design courses, I usually organize learning around clear design tasks:
\begin{itemize}
  \item analysing user needs and environmental constraints before form-making;
  \item connecting spatial layout, material choice, lighting, circulation, safety, accessibility, and comfort;
  \item using iterative critiques so students can improve their design decisions through evidence and feedback;
  \item asking students to explain design decisions verbally, visually, and technically.
\end{itemize}

This approach is especially important for age-friendly and inclusive design. Students need to understand that designing for older adults is not a decorative theme. It requires knowledge of mobility, visibility, cognition, comfort, safety, dignity, social participation, and care practice. I therefore integrate user-centred and evidence-based thinking into studio assignments whenever possible.

\section*{Teaching Decorative Materials and Technical Courses}
For Decorative Materials, CAD/BIM, design software, building structures, and renovation-related courses, I focus on the link between technical knowledge and design judgement. Students should not only memorize materials or software commands; they should understand why a material, construction method, drawing standard, or digital workflow matters in real design practice.

My teaching strategy includes:
\begin{itemize}
  \item using material samples, case studies, and construction details to connect theory with application;
  \item teaching CAD/BIM/software skills through design tasks rather than isolated commands;
  \item using precedent analysis to show how material, structure, cost, maintenance, safety, and sustainability influence design outcomes;
  \item encouraging students to compare alternatives and justify decisions through criteria such as accessibility, durability, health, environmental performance, and user comfort.
\end{itemize}

This is useful for Malaysian built environment education because students need both creative design capacity and employable technical competence. I can contribute to courses that connect interior design, architecture, building technology, renovation, digital representation, and professional practice.

\section*{Integrating AI and Research into Teaching}
My research on AI-assisted age-friendly environment assessment can enrich teaching in several ways. I would introduce students to responsible uses of AI in design education: extracting design precedents, analysing user needs, generating early design alternatives, checking accessibility criteria, evaluating environmental risks, and communicating design evidence. However, I would also teach students that AI outputs must be critically reviewed, ethically used, and grounded in professional judgement.

In advanced undergraduate or postgraduate teaching, I can develop research-led projects such as:
\begin{itemize}
  \item AI-assisted assessment of age-friendly residential interiors;
  \item BIM-linked accessibility evaluation for housing or community facilities;
  \item material and spatial strategies for elderly care environments;
  \item renovation proposals for existing buildings serving older adults;
  \item evidence-based design guidelines for inclusive and healthy environments.
\end{itemize}

These topics can produce strong studio outcomes, student research projects, and potential publications or community engagement outputs.

\section*{Assessment and Student Development}
I believe assessment should reward process, evidence, technical clarity, creativity, and professional communication. For design courses, I would use a balanced assessment structure that includes research/context analysis, concept development, interim critique, technical drawings, final presentation, and reflective writing. For technical courses, I would combine exercises, applied assignments, quizzes, and project-based submissions.

I also value formative feedback. Many design students improve significantly when feedback is specific and staged. I therefore prefer to give comments that identify the design problem, explain the reason, and suggest the next action. This helps students learn how to improve independently rather than simply follow instructions.

\section*{Contribution to Malaysian University Teaching}
For Malaysian lecturer positions, I can contribute to modules in Interior Design, Interior Architecture, Architecture, Built Environment, Decorative Materials, CAD/BIM, design software, renovation design, universal design, and age-friendly design. I can teach undergraduate students and support studio/project supervision. Over time, I would also like to develop elective or project-based teaching on AI for inclusive built environment design, connecting digital methods with Malaysia's ageing society and built environment needs.

\end{document}
